\section{真题速看}
\subsection{打酱油}
\textbf{第11次CCF-CSP认证，2017-09，第1题。}\href{https://www.acwing.com/problem/content/3250/}{打酱油（在线做题）}。
\textbf{题目与输入输出。}带N元买每瓶10元的酱油，可选择买3瓶赠1瓶或买5瓶赠2瓶，促销组之间不重复使用已赠送的瓶。求最多得到多少瓶。输入N，输出最大瓶数。N是10的整数倍且不超过300。
\textbf{读题思路。}先把钱换成可付费的瓶数，赠品不消耗钱，也不能把同一份购买重复参加两种促销。付费瓶数至多30，枚举促销组完全够用。
\textbf{做题思路与解答。}枚举买5赠2的组数f。剩余付费瓶数为r时，尽可能组成买3赠1的组，剩下的按单瓶购买，得到$7f+4\lfloor r/3\rfloor+r\bmod3$。每种方案都有一个f；固定f时多组成三瓶组只会增加赠品，所以逐项取最大值必然包含最优方案。也可证明优先五瓶组的贪心，但本题的小规模枚举不需要额外假设。
\begin{verbatim}
输入
40
输出
5
\end{verbatim}
\begin{verbatim}
输入
80
输出
11
\end{verbatim}
\textbf{复杂度与边界。}时间$O(N/50+1)$，额外空间$O(1)$。0元时答案0；余下4个付费瓶应组成三瓶组加单瓶。
\Code{49_soy.cpp}
\subsection{跳一跳}
\textbf{第13次CCF-CSP认证，2018-03，第1题。}\href{https://www.acwing.com/problem/content/3260/}{跳一跳（在线做题）}。
\textbf{题目与输入输出。}输入若干个0、1、2，数量不超过30；0仅在末尾出现，表示游戏结束。1表示跳到方块非中心，得1分；2表示跳到中心，若是第一次跳或上次得1分，则得2分，否则比上次多2分。输出总得分。
\textbf{读题思路。}输入没有n，由0终止；连续中心得分为2、4、6等，遇到1后必须重置。末尾0不计分。
\textbf{做题思路与解答。}last表示上一跳实际得分，初值0。读到1令last=1；读到2且last不超过1时令last=2，否则加2。然后累加last。每次更新恰好实现规则，因此处理完前缀后总分始终正确。
\begin{verbatim}
输入
1 1 2 2 2 1 1 2 2 0
输出
22
\end{verbatim}
\textbf{复杂度与边界。}时间与输入数量成正比，空间$O(1)$。输入首项为0时总分为0；不要把last当作上一次输入的1或2。
\Code{50_jump.cpp}
\subsection{线性分类器}
\textbf{第19次CCF-CSP认证，2020-06，第1题。}\csp{CSP202006A}{线性分类器}。
\textbf{题目与输入输出。}给出n个点，每点含整数坐标x、y以及类别A或B，两类都存在。随后给m条直线$\theta_0+\theta_1x+\theta_2y=0$；保证没有给定点恰在直线上。若同类点全部位于同一侧、异类点位于另一侧，输出Yes，否则No，每条直线一行。输入首行n、m，随后n行点、m行三个系数。$1\le n\le1000$，$1\le m\le20$，坐标和系数绝对值不超过$10^4$。
\textbf{评测限制。}曙梦OJ当前时限1秒，内存512MB。
\textbf{读题思路。}必须把一条线与所有点比较；不是检查两类各找一个代表就结束。Yes和No的大小写要精确。两侧由表达式符号区分，无需求斜率，因而竖直线也适用。
\textbf{做题思路与解答。}用第一个点确定“正侧应该是哪一类”，再检查每个点的符号是否与类别相符。如果有反例，这条线不满足。代入式使用long long计算；原范围下int也能容纳，但显式提升类型便于保持安全的乘法习惯。正确性来自直线两侧的表达式符号相反，以及逐点检查完整覆盖了分类条件。
\begin{verbatim}
输入
9 3
1 1 A
1 0 A
1 -1 A
2 2 B
2 3 B
0 1 A
3 1 B
1 3 B
2 0 A
0 2 -3
-3 0 2
-3 1 1
输出
No
No
Yes
\end{verbatim}
\textbf{复杂度与边界。}时间$O(nm)$，空间$O(n)$。结构体记录三项属性；读类别用格式串中的空格跳过之前的换行。
\Code{38_classifier.cpp}
\subsection{检测点查询}
\textbf{第20次CCF-CSP认证，2020-09，第1题。}\href{https://www.acwing.com/problem/content/3295/}{检测点查询（在线做题）}。
\textbf{题目与输入输出。}输入n、居民坐标X、Y，随后n行检测点坐标，编号按输入顺序从1开始。输出最近的三个点编号，每个一行；同距离时编号小者在前。$3\le n\le200$，坐标绝对值不超过1000。
\textbf{读题思路。}要按距离、编号两层比较。无需真的开方，平方距离与距离具有相同顺序。只取三个点，三轮扫描就够用。
\textbf{做题思路与解答。}计算各点平方距离。每轮从未选点中找最小距离者，按编号升序扫描且只在距离严格减小时替换，就自动保留同距离最小编号。输出并标记已选，重复三轮。归纳可知每轮得到剩余点中的第一名，三轮恰好得到所需前三名。
\begin{verbatim}
输入
3 2 2
2 2
2 3
2 4
输出
1
2
3
\end{verbatim}
\begin{verbatim}
输入
5 0 1
-1 0
0 0
1 0
0 2
-1 2
输出
2
4
1
\end{verbatim}
\textbf{复杂度与边界。}时间$O(3n)=O(n)$，数组空间$O(n)$。平方距离最大为$8\times10^6$，此范围int可装下；代码用long long并先提升差值类型。
\Code{51_nearest.cpp}
\subsection{田地丈量}
\textbf{第29次CCF-CSP认证，2023-03，第1题。}\csp{CSP202303A}{田地丈量}。
\textbf{题目与输入输出。}输入n、a、b，随后n行现有田地矩形的左下与右上坐标。现有田地互不重叠（交集面积为0），目标区域为$[0,a]\times[0,b]$。求目标区域内已有田地面积。$n\le100$，所有坐标绝对值不超过$10^4$，a、b为正。输出面积。
\textbf{读题思路。}统计连续面积，边接触贡献0，不能加1。现有田地互不重叠，所以各块与目标的交面积可以直接相加；若取消这一保证，则不能简单累加。
\textbf{做题思路与解答。}每块田地裁剪到目标：横向左端取$\max(x_1,0)$、右端取$\min(x_2,a)$，纵向同理。有正宽、正高才累加其乘积。两个一维交区间的笛卡尔积正好是矩形交集；原田地互不重叠，因此求和无重复覆盖。
\begin{verbatim}
输入
4 10 10
0 0 5 5
5 -2 15 3
8 8 15 15
-2 10 3 15
输出
44
\end{verbatim}
\textbf{复杂度与边界。}时间$O(n)$，额外空间$O(1)$，可边读边算。面积与累加值用long long，乘法前提升类型。
\Code{52_land.cpp}
